\documentclass[prb,preprint]{revtex4-1}
% AJP uses the same style as Phys. Rev. B (prb).

\usepackage{amsmath}
\usepackage{amsfonts}

\begin{document}

\title{Deriving quantum mechanics from probabilistic classical mechanics}

% Anonymous for review submission:
%\author{Warren MacEvoy}
%\affiliation{Department of Electrical and Computer Engineering, Colorado Mesa University, Grand Junction, CO 81501}

\date{\today}

\begin{abstract}
We show that non-relativistic quantum mechanics can be derived from classical
Hamiltonian dynamics with probabilistic initial data by lifting the equations of motion
to a Liouville equation, Fourier transforming in momentum, and applying a centered-difference
approximation.  The resulting equation separates into forward and backward
Schr\"odinger equations.  The only approximation is $O(\hbar^3)$, rendering it
indistinguishable from exact.  The Born rule emerges as a time-averaged marginal
over momentum.  We discuss implications for the uncertainty principle, entanglement,
and computation.
\end{abstract}

\maketitle


\section{Introduction}

The goal of these notes is to show non-relativistic quantum mechanics can be closely
approximated by classical dynamics with probabilistic initial data, along with an interesting
transformation of coordinates.  The approximation may be so close as to be indistinguishable.
We give the notes in terms of one spatial dimension, but it is trivial to extend to higher
dimensions and multiple particles.


\section{Quantum Perspective}

Consider a tiny one-dimensional non-relativistic particle of mass $M$ in a tiny potential $V(x)$.
Here the accepted model is the Schr\"odinger equation, where $|\Psi(x,t)|^2$ gives the probability
of observing the particle at location $x$ at time $t$.  The wave function $\Psi(x,t)$
evolves according to the partial differential equation
\begin{equation}\label{eq:schrodinger}
i\hbar \Psi_t = -\frac{\hbar^2}{2M}\Psi_{xx} + V(x)\Psi \,,
\end{equation}
where $\hbar \approx 1.05 \times 10^{-34}$~joule-seconds.


\section{Classical Perspective}

Consider a regular one-dimensional non-relativistic particle of mass $M$ in a regular potential $V(x)$.
The position of the particle $x(t)$ and momentum of the particle $p(t)$ evolve according
to the ordinary differential equation
\begin{eqnarray}
M \frac{dx}{dt} & = & p \,, \label{eq:hamilton-x}\\
\frac{dp}{dt} & = & -\frac{\partial V}{\partial x} \,. \label{eq:hamilton-p}
\end{eqnarray}


\section{Classical from Quantum Perspective}

There should be a way to see classical mechanics arrive as an approximation of quantum
mechanics.  This should not be surprising, and is included here for completeness.\cite{landau1977}

Approximate
\begin{equation}
\Psi \approx A \exp\left(\frac{iS}{\hbar}\right) \,,
\end{equation}
and, keeping the highest order terms:
\begin{equation}
-S_t = \frac{1}{2M}\left(S_x\right)^2 + V \,.
\end{equation}
Differentiate with respect to $x$:
\begin{equation}
-S_{xt} = \frac{1}{M}S_x S_{xx} + \frac{\partial V}{\partial x} \,.
\end{equation}
Now define $P = S_x$
\begin{equation}
-P_t = \frac{1}{M}P P_x + \frac{\partial V}{\partial x} \,.
\end{equation}
This is a first order hyperbolic differential equation and is solvable by the method of
characteristics.  Using $p(t) = P(x(t),t)$:
\begin{equation}
\frac{dp}{dt} = \frac{\partial P}{\partial x}\frac{dx}{dt} + \frac{\partial P}{\partial t}
= \frac{\partial P}{\partial x}\frac{dx}{dt} - \frac{1}{M}P P_x - \frac{\partial V}{\partial x} \,.
\end{equation}
Choosing
\begin{eqnarray}
\frac{dx}{dt} &=& \frac{1}{M}p \,, \\
\frac{dp}{dt} &=& -\frac{\partial V}{\partial x}
\end{eqnarray}
reduces the equation.

So, to leading order, the Schr\"odinger equation carries phase-gradient
information along characteristics defined by the equations of motion for classical mechanics.
Note that characteristics will inevitably be singular.


\section{Three concepts to get quantum mechanics from classical mechanics}

So what has to be added to classical mechanics to get quantum mechanics?  This seems more
interesting.  You could reverse the WKB approximation above, but that is not particularly
enlightening.  We instead combine the following concepts.

\subsection{An ordinary differential equation can be lifted to a linear partial differential equation}

Think of an ordinary differential equation initial value problem as evolving a
partial differential equation which describes the evolution of the probability
distribution of the initial values.\cite{goldstein2002}

ODE:
\begin{eqnarray}
\frac{dx_1}{dt} &=& v_1 \,, \\
\frac{dx_2}{dt} &=& v_2 \,, \\
&\vdots& \nonumber \\
\frac{dx_n}{dt} &=& v_n \,.
\end{eqnarray}

PDE:
\begin{equation}\label{eq:liouville}
\frac{\partial f}{\partial t}
+ \frac{\partial (v_1 f)}{\partial x_1}
+ \frac{\partial (v_2 f)}{\partial x_2}
+ \cdots
+ \frac{\partial (v_n f)}{\partial x_n}
= 0 \,.
\end{equation}

Using the Dirac delta function for the initial value, the PDE recovers the ODE initial value
problem: $f(x,0) = f_0(x) = \delta(x-x_0)$.

\subsection{The Fourier transform}

\begin{equation}\label{eq:fourier-forward}
\hat{\Psi}(\hat{p}) = \int_{-\infty}^{+\infty} dp \, e^{-i\hat{p}p/\hbar}\,\Psi(p) \,,
\end{equation}
\begin{equation}\label{eq:fourier-inverse}
\Psi(p) = \int_{-\infty}^{+\infty} \frac{d\hat{p}}{2\pi \hbar} \, e^{+i\hat{p}p/\hbar}\,\hat{\Psi}(\hat{p}) \,,
\end{equation}
\begin{equation}
\partial_p \Psi \rightarrow i \frac{\hat{p}}{\hbar} \hat{\Psi} \,,
\end{equation}
\begin{equation}
p\Psi \rightarrow i \hbar \partial_{\hat{p}} \hat{\Psi} \,.
\end{equation}
\begin{equation}\label{eq:delta}
\delta(\hat{p})= \int_{-\infty}^{+\infty} \frac{dp}{2\pi \hbar} \, e^{i\hat{p}p/\hbar} \,.
\end{equation}
The use of $\hat{p} = \hbar \omega$ compared to the typical spectral parameter $\omega$ greatly
compresses the spectrum.

\subsection{An approximation of $V'$ by a centered difference}

\begin{equation}\label{eq:centered-diff}
\frac{\partial V(x)}{\partial x} \approx \frac{V(x+\varepsilon/2) - V(x-\varepsilon/2)}{\varepsilon}
+ O(\varepsilon^2)  \,.
\end{equation}
We actually approximate $\varepsilon V'$, so with first error term:
\begin{equation}\label{eq:centered-diff-error}
\varepsilon \frac{\partial V(x)}{\partial x} \approx V(x+\varepsilon/2) - V(x-\varepsilon/2)
+ \frac{\varepsilon^3}{24} \frac{\partial^3 V(x)}{\partial x^3} + O(\varepsilon^5)  \,.
\end{equation}


\section{Quantum from probabilistic classical mechanics}

We note the structure here is related to the Wigner-Moyal formulation,\cite{wigner1932,moyal1949,hillery1984}
arrived at here from the classical rather than quantum direction.

Using the ODE $\Rightarrow$ PDE concept for the classical Hamiltonian system, the ODE is
\begin{eqnarray}
\frac{dx}{dt} &=& \frac{1}{M}p \,, \\
\frac{dp}{dt} &=& -\frac{\partial V}{\partial x} \,.
\end{eqnarray}

This corresponds to the PDE describing the evolution of the probability density:
\begin{equation}
\frac{\partial f}{\partial t}
+ \frac{\partial}{\partial x}\left(\frac{1}{M}pf\right)
+ \frac{\partial}{\partial p}\left(-\frac{\partial V}{\partial x}f\right)
= 0 \,,
\end{equation}
or
\begin{equation}\label{eq:liouville-hamiltonian}
f_t + \frac{p}{M} f_x - \frac{\partial V}{\partial x} f_p = 0 \,.
\end{equation}

Now Fourier transform only with respect to the momentum $p$.  This results in
\begin{equation}\label{eq:fourier-liouville}
i\hbar \hat{f}_t - \frac{\hbar^2}{M} \hat{f}_{x\hat{p}} + \hat{p} \frac{\partial V}{\partial x} \hat{f} = 0 \,.
\end{equation}

We now change variables and use the only approximation.  Define
\begin{equation}
z^{\pm} = x \pm \frac{\hat{p}}{2} \,,
\end{equation}
so
\begin{eqnarray}
\frac{\partial}{\partial x} & = & \frac{\partial}{\partial z^+} + \frac{\partial}{\partial z^-}\,,\\
\frac{\partial}{\partial \hat{p}} &=& \frac{1}{2}\frac{\partial}{\partial z^+}
- \frac{1}{2}\frac{\partial}{\partial z^-}
\end{eqnarray}
and use the approximation~(\ref{eq:centered-diff}) with $\varepsilon = \hat{p}$:
\begin{equation}\label{eq:approx-V}
\frac{\partial V}{\partial x}
\approx \frac{V(x+\hat{p}/2) - V(x-\hat{p}/2)}{\hat{p}} \,.
\end{equation}

This transforms Eq.~(\ref{eq:fourier-liouville}) into
\begin{equation}\label{eq:f-hat-pde}
i\hbar \hat{f}_t - \frac{\hbar^2}{2M}
\left[\left(\frac{\partial}{\partial z^{+}}\right)^2 - \left(\frac{\partial}{\partial z^{-}}\right)^2\right]\hat{f}
+ [V(z^+) - V(z^-)]\hat{f} = 0 \,.
\end{equation}
While this does not yet look like Schr\"odinger's equation, this is a linear differential
equation, and we can look at separation of variables to construct a general solution
to Eq.~(\ref{eq:f-hat-pde}).
Suppose $\hat{f} = \sum \psi^+(z^+)\psi^-(z^-)$ where each of these terms satisfy the above; we have
\begin{equation}\label{eq:separated}
\pm i\hbar \psi_t^{\pm} = -\frac{\hbar^2}{2M}\psi^{\pm}_{zz} + V\psi^{\pm} \,.
\end{equation}
Here $z$ takes the respective roles of $z^+$ and $z^-$.  For real potential $V$, we can have
$\psi^+ = e^{i\lambda t} \psi_\lambda$ and $\psi^- = e^{-i\lambda t}\psi_\lambda^*$ where the $\psi_\lambda$ are eigenvectors
of the original Schr\"odinger equation~(\ref{eq:schrodinger}).  These form a complete basis for the
solution of PDE for $\hat{f}$:
\begin{equation}\label{eq:f-hat-expansion}
\hat{f}(t,z_+,z_-)=\sum_{\lambda,\lambda'} C_{\lambda,\lambda'} e^{i(\lambda-\lambda')t} \psi_{\lambda}(z_+) \psi^*_{\lambda'}(z_-) \,,
\end{equation}
where the $C_{\lambda,\lambda'}$ are determined by the initial distribution $f_0$.
Assuming $\psi_\lambda$ are normalized by $\int_{-\infty}^{+\infty} dz \, |\psi_\lambda(z)|^2 = 1$, we have:
\begin{eqnarray}
C_{\lambda,\lambda'} &=& \int_{-\infty}^{+\infty} dz^+ \int_{-\infty}^{+\infty} dz^- \, \psi_\lambda^*(z^+) \psi_{\lambda'}(z^-) \hat{f}_0(z^+,z^-) \,, \label{eq:C-z}\\
&=& \int_{-\infty}^{+\infty} dx \int_{-\infty}^{+\infty} d\hat{p} \, \psi_\lambda^*\!\left(x+\frac{\hat{p}}{2}\right) \psi_{\lambda'}\!\left(x-\frac{\hat{p}}{2}\right) \hat{f}_0(x,\hat{p}) \,. \label{eq:C-xp}
\end{eqnarray}


\section{The Born projection}

Now let us compute the probability density function $\rho(x)$ of observed position by integrating $f$ over a time window and the momentum:
\begin{eqnarray}
\rho(x) &=& \lim_{T\to\infty} \frac{1}{T} \int_0^T dt \int_{-\infty}^{+\infty} dp \, f(x,p) \,, \label{eq:born-1}\\
&=& \lim_{T\to\infty} \frac{1}{T} \int_0^T dt \int_{-\infty}^{+\infty} dp \int_{-\infty}^{+\infty} \frac{d\hat{p}}{2\pi \hbar} \, e^{i\hat{p}p/\hbar} \hat{f} \,, \label{eq:born-2}\\
&=& \lim_{T\to\infty} \frac{1}{T} \int_0^T dt \int_{-\infty}^{+\infty} d\hat{p} \, \hat{f} \int_{-\infty}^{+\infty} \frac{dp}{2\pi \hbar} \, e^{i\hat{p}p/\hbar} \,, \label{eq:born-3}\\
&=& \lim_{T\to\infty} \frac{1}{T} \int_0^T dt \int_{-\infty}^{+\infty} d\hat{p} \, \hat{f}\,\delta(\hat{p}) \,, \label{eq:born-4}\\
&=& \lim_{T\to\infty} \frac{1}{T} \int_0^T dt \, \hat{f}_{\hat{p}=0} \,, \label{eq:born-5}\\
&=& \lim_{T\to\infty} \frac{1}{T} \int_0^T dt \sum_{\lambda,\lambda'} C_{\lambda,\lambda'} e^{i (\lambda-\lambda')t}\psi_\lambda(x)\psi_{\lambda'}^*(x) \,, \label{eq:born-6}\\
&=& \sum_\lambda C_{\lambda,\lambda} |\psi_\lambda(x)|^2 \,. \label{eq:born-rule}
\end{eqnarray}
Notice it does not affect the dynamics---it is merely the observable projection of the dynamics
to a measurement.  We call this the Born projection, which is the Born rule, but now derived
and without decoherence.\cite{zurek2003}  This happens to be coupled to decoherence with any practical
measurement device we have so far constructed.


\section{Demystifications}

\subsection{Heisenberg's uncertainty principle}

There is nothing wrong with a distribution that specifies the precise initial position and
momentum of the particle.  However, we can only observe the Born projection---which makes
the view fuzzy.

Substitution of $f_0(x,p)=\delta(x-x_0)\delta(p-p_0)$ into the formula for $C_{\lambda,\lambda'}$
gives the cross-Wigner function as the coefficients of the expansion of $\hat{f}$ in the
eigenbasis of the Schr\"odinger equation:
\begin{equation}\label{eq:cross-wigner}
C_{\lambda,\lambda'} = \int_{-\infty}^{+\infty} d\hat{p} \, \psi_\lambda^*\!\left(x_0+\frac{\hat{p}}{2}\right) \psi_{\lambda'}\!\left(x_0-\frac{\hat{p}}{2}\right) e^{-i\hat{p}p_0/\hbar}\,.
\end{equation}
The Born projection makes this look imprecise.

\subsection{Quantum entanglement}

The above derivation is for a single particle in a single space dimension, but it is trivial
to extend to multiple particles in three spatial dimensions.  Entanglement amounts to a
dependent choice of initial distribution, so $f$ cannot be factored into independent components.
The underlying mechanism evolves independently; it is only the initial distribution and final
Born projection that makes it look like there is something mysterious.  It is actually a simple
consequence of dependent distributions and the Born projection.

For two particles and some single-particle eigenstate $\psi_\lambda(x)$, choosing
\begin{equation}\label{eq:entanglement}
\hat{f}_0=\frac{1}{\sqrt{2}}\psi_\lambda\!\left(x_1+\tfrac{1}{2}\hat{p}_1\right)\psi^*_\lambda\!\left(x_1-\tfrac{1}{2} \hat{p}_1\right) +\frac{1}{\sqrt{2}}\psi_\lambda\!\left(x_2+\tfrac{1}{2}\hat{p}_2\right)\psi^*_\lambda\!\left(x_2-\tfrac{1}{2} \hat{p}_2\right) \,,
\end{equation}
where $x_{1,2}$ and $p_{1,2}$ are the 3-dimensional position and momentum coordinates of the two
particles and $f$ is a probability distribution in $\mathbb{R}^{12}$.  This is the standard EPR entanglement
in this representation.  The Born projection then gives the observed correlation.

\subsection{Quantum computation}

The power of a quantum computer is only in the Fourier transform.  If the machine setup
starts from a fixed wave function, and therefore fixed distribution $f$, then the machine can
only apply classical dynamics to the evolution of that distribution.  The only power is the
Born projection, which computationally is not interesting except for the Fourier transform.

Even if you prefer the standard form of quantum mechanics, this statement is still true---the
formulations are indistinguishable.


\section{Discussion}

The only approximation here is the centered difference estimate:
\begin{equation}\label{eq:approximation}
\hat{p} \frac{\partial V}{\partial x}
\approx V\!\left(x+\frac{\hat{p}}{2}\right) - V\!\left(x-\frac{\hat{p}}{2}\right)
= \hat{p} \frac{\partial V}{\partial x} + \frac{\hat{p}^3}{24}\frac{\partial^3 V}{\partial x^3} + \cdots \,.
\end{equation}
The Fourier transform~(\ref{eq:fourier-forward}) compresses the spectrum of momentum by a factor
of $\hbar$, so the approximation~(\ref{eq:approximation}) is $O(\hbar^3)$, or about $10^{-100}$ of
human scales.  This is indistinguishable from exact for all practical purposes.  Even for
results where $V$ is not smooth, an $\hbar^2$-scale smoothing of the potential is indistinguishable
from the original, so the approximation is effectively exact for all observational purposes.

Regarding Bell's theorem,\cite{bell1964,bell1987} note that the introduced hidden variables (classical
momentum) are natural from the classical perspective, and the Fourier transform is the
non-locality which allows this reformulation to co-exist with Bell's theorem.  The mathematics
is indifferent to the interpretation; they are both effectively equivalent, but the interpretation
is compelling from the simplicity of the underlying mechanism.

These notes stayed in a drawer for a long time, since I was not convinced they explained
anything.  But they provide a useful demystification of quantum mechanics, allowing the intuition
of classical mechanics to provide insight.  They place the mystery specifically: why does the
universe conspire to transform the momentum space, and choose to reveal only the DC terms to
us as observers?  No good answer to this kept these notes in my drawer, but others may have
better ideas.  They also suggest a global unknowable state dependence through the transform
that survives observation.  Decoherence is real, but maybe a side-effect of large interactions
but not as central as postulated.

I hope others find these ideas interesting.


\begin{acknowledgments}
The author has no conflicts to disclose.
\end{acknowledgments}


\begin{thebibliography}{9}

\bibitem{landau1977} L. D. Landau and E. M. Lifshitz, \textit{Quantum Mechanics:
Non-Relativistic Theory}, 3rd edition (Pergamon Press, Oxford, 1977).

\bibitem{goldstein2002} Herbert Goldstein, Charles Poole, and John Safko,
\textit{Classical Mechanics}, 3rd edition (Addison-Wesley, San Francisco, 2002).

\bibitem{wigner1932} E. P. Wigner, ``On the quantum correction for thermodynamic
equilibrium,'' Phys. Rev. \textbf{40}, 749--759 (1932).

\bibitem{moyal1949} J. E. Moyal, ``Quantum mechanics as a statistical theory,''
Math. Proc. Cambridge Philos. Soc. \textbf{45}, 99--124 (1949).

\bibitem{hillery1984} M. Hillery, R. F. O'Connell, M. O. Scully, and E. P. Wigner,
``Distribution functions in physics: Fundamentals,''
Phys. Rep. \textbf{106} (3), 121--167 (1984).

\bibitem{zurek2003} W. H. Zurek, ``Decoherence, einselection, and the quantum
origins of the classical,'' Rev. Mod. Phys. \textbf{75} (3), 715--775 (2003).

\bibitem{bell1964} J. S. Bell, ``On the Einstein Podolsky Rosen paradox,''
Physics \textbf{1} (3), 195--200 (1964).

\bibitem{bell1987} J. S. Bell, \textit{Speakable and Unspeakable in Quantum Mechanics}
(Cambridge University Press, Cambridge, 1987).

\end{thebibliography}

\end{document}
